\section{合同的权利义务终止}
\begin{enumerate}
  \item \zd{清}偿。考清偿不了。
        \begin{enumerate}
          \item 代物清偿协议/以物抵债（新债）。\zd{借款期限届满}（原债）。\al{用建设用地使用权抵偿 2000W 债务}。
                \begin{enumerate}
                  \item \zd{诺成}（最高法观点）。自由选择物、钱（原债并未消灭）。\al{只能以交付古画履行义务 ×}。
                  \item 物有瑕疵 \yjt 瑕疵担保责任 \yjt 违约责任。\al{顾华交付后如果被鉴定为赝品}。
                \end{enumerate}
          \item 清偿抵充。欠多笔钱还的是哪一笔，排序问题。
                \begin{enumerate}
                  \item 约定 > 债务人指定 > \zd{到期的 > 缺乏担保/少} > 负担重（利息高）> 到期先后顺序 > 按比例。
                \end{enumerate}
          \item 清偿顺序。有关费用 > 利息 > 主债务。
        \end{enumerate}
  \item 合同\zd{解}除。未履行完毕 + \zd{合法有效} + 根本违约（目的无法实现）。必考。\al{双方交割完毕}。
        \begin{enumerate}
          \item 类型。
                \begin{enumerate}
                  \item 双方解除/协议解除/合意解除。\zd{千万千万别看内容，看主体，还是容易掉坑里}！
                        \al{甲乙签订协议 1，乙丙丁签订协议 2，约定协议 2 签订之日协议 1 自动解除。协议 1 被协议 2 取代 ×}。
                  \item 单方解除。\yi 约定解除。约定的解除情况出现。\er 法定解除。\al{\xhx{死亡}不一定是合同解除的法定事由，切记}。
                        \begin{enumerate}
                          \item 一般法定解除。可能存在\zd{权利竞合}。违约责任 \hz 缔约过失责任。
                                \begin{itemize}
                                  \item 预期违约、迟延履行、\zd{从给付义务的违反}（\al{买房养老设施}）\yjt 先催告。
                                  \item 不可抗力。不能预见、不可避免、不能克服的客观情况。\zd{唯一法定免责事由}。
                                  \item 情势变更。继续履行成本太高，显失公平 \yjt 协商前置。\al{国家政策、因地震导致原材料成本和运输费用上涨}。
                                \end{itemize}
                          \item 特殊法定解除/任意解除。\zd{人身、服务类、无期限}相关。合理期限通知。1 个选项，问谁有。
                                \begin{itemize}
                                  \item \zd{无期限则无逾期}。
                                  \item 定作人、托运人、不定期租赁单方；不定期保管、委托合同双方。
                                \end{itemize}
                        \end{enumerate}
                \end{enumerate}
          \item 行使。通知解除、行为解除（或诉或裁）。
          \item \zd{期限。可约定}。
          \item 法律效果。\yi \zd{不影响合同效力、违约责任、担保责任}。\er 补救措施 \yjt 恢复原状。
          \item vs 撤销权。（典型的）形成权。
        \end{enumerate}
  \item \zd{抵}销。互负债权债务。或诉或裁或通知。
        \begin{enumerate}
          \item \zd{法定抵销}。
                \begin{enumerate}
                  \item 互债、同类。考试都是货币之债。
                  \item \zd{主动债权}已到期。谁提抵销谁就是主动方。\al{甲欠乙 10W，12/20 到期，乙欠甲 5W，12/21 到期，甲提抵销}。
                  \item \zd{到期在先原则}。谁先到期保护谁。
                  \item \zd{已过诉讼时效之债的债权，不得作为主动债权来主张抵销}。
                  \item 数额只看出质人之间……\al{甲将乙的 10W 债权出质给银行，担保其 9W 的贷款，出质前，乙对甲享有2万元到期债权。银行可行使的最大数额为 8}。
                \end{enumerate}
          \item 意定抵销。协商一致，意思自治。没啥考点。
        \end{enumerate}
  \item \zd{提}存。\zd{债权人无正当理由拒绝受领债权}。\zd{清偿（债权效力）、担保（物权效力）}、执行、保管。
        \begin{enumerate}
          \item 视为履行，债务消灭。
          \item 提存部门妥善保管义务。
          \item 提存人凭判决、裁定、公证证明取回提存物。
        \end{enumerate}
  \item \zd{免}除。未履行完毕
  \item \begin{enumerate}
          \item 单方行为说，类推动产抛弃。
          \item \zd{修正的单方行为说}。债务人合理期限明确反对除外。\al{甲对乙表示，丙生活困难，这钱他不用还了}。
        \end{enumerate}
  \item \zd{混}同。\zd{结婚不会发生混同，叫免除}。\al{甲公司被乙公司收购}。
\end{enumerate}
